\section{Gated networks}\label{sec:networks}

We now set up the class of networks and the reduction that turns a differentiable network into a
product of gated Gaussian blocks. The reduction is exact for piecewise-linear activations and holds
in a convex-hull sense for any activation with derivative in $[0,1]$; this is what lets a single set
of estimates cover the rectified linear unit, its leaky and absolute-value variants, and smooth
activations at once.

\subsection{Feedforward networks with biases}
A depth-$L$ fully connected network is a map $F:\R^{d}\to\R$,
\begin{equation}\label{eq:fcnet}
  h_0=x,\qquad h_\ell=\sigma_\ell\!\bigl(W_\ell h_{\ell-1}+b_\ell\bigr)\ (1\le\ell\le L),\qquad
  F(x)=\langle w, h_L\rangle,
\end{equation}
with weight matrices $W_\ell\in\R^{n_\ell\times n_{\ell-1}}$ ($n_0=d$), bias vectors
$b_\ell\in\R^{n_\ell}$, a readout $w\in\R^{n_L}$, and activations $\sigma_\ell$ applied
coordinatewise. We take the weights independent Gaussian, $W_\ell$ with entry variance
$\theta_\ell/n_{\ell-1}$ and $w$ with variance $1/n_L$, and treat the biases as arbitrary fixed vectors,
subject only to a mild condition below that keeps the gates from saturating. The scale $\theta_\ell$ is
the \emph{signal-preserving} choice ($\theta_\ell=2$ for ReLU, so that $\EE\|h_\ell\|^2=\|h_{\ell-1}\|^2$)
where the gradient lower bound of \Cref{prop:gradlb} needs it, and $\theta_\ell=1$ elsewhere; all
constants below are stated for the choice in force. Binary classification is by $\sgn F$; the
multiclass case has a readout matrix and $F(x)=\arg\max_j\langle w_j,h_L\rangle$, treated in
\Cref{sec:fragility}.

\subsection{The derivative gate}
Wherever $F$ is differentiable, the chain rule gives
\begin{equation}\label{eq:jacobian}
  \nabla F(x)^\top=w^\top D_L(x)W_L\,D_{L-1}(x)W_{L-1}\cdots D_1(x)W_1,\qquad
  D_\ell(x)=\diag\bigl(\sigma_\ell'(z_\ell(x))\bigr),
\end{equation}
where $z_\ell(x)=W_\ell h_{\ell-1}(x)+b_\ell$ is the preactivation. The diagonal $D_\ell(x)$ is the
\emph{gate} at layer $\ell$. Three cases organize the activations.

\begin{definition}[Gate class]\label{def:gate}
An activation $\sigma$ is \emph{$[0,1]$-slope} if it is Lipschitz with $\sigma'\in[0,1]$ almost
everywhere. It is \emph{piecewise linear} if in addition $\sigma'$ takes finitely many values
$s_1,\dots,s_r\in[0,1]$.
\end{definition}

For a piecewise-linear activation the gate $D_\ell(x)$ is, coordinatewise, one of the finitely many
slopes, so $D_\ell(x)$ ranges over a finite set of diagonal matrices as $x$ varies; for the rectified
linear unit the slopes are $\{0,1\}$ and $D_\ell(x)$ is a coordinate projection. For a general
$[0,1]$-slope activation the gate is a diagonal matrix with entries in $[0,1]$, hence a convex
combination of $0/1$ gates. Since the operator norm and least singular value of $DW$, as $D$ ranges
over diagonal matrices with entries in $[0,1]$ and fixed support pattern, are attained at extreme
points (the map $D\mapsto\|DWx\|$ is convex, and $D\mapsto\smin(DW)$ is concave on the order interval),
every bound we prove for $0/1$ gates transfers to the convex hull. We therefore state the estimates
for $0/1$ gates and use them verbatim for any $[0,1]$-slope activation; this is recorded as
\Cref{lem:hull} in \Cref{app:netlemmas}.

\subsection{Realized patterns and gate density}
For an input $x$, the \emph{realized pattern} at layer $\ell$ is the support of $D_\ell(x)$ (the set of
neurons with nonzero slope), of size $|{\supp D_\ell(x)}|=:\alpha_\ell(x)\,n_\ell$. Under a symmetric
input law and an odd activation the preactivations are symmetric and each gate has density $1/2$ in
the limit; with a nonzero bias $b_\ell$ the density is set by the bias level, and we assume it stays
bounded away from $0$ and $1$.

\begin{assumption}[Nondegenerate gates]\label{ass:nondeg}
There are constants $0<c_0\le c_1<1$ such that, w.h.p.\ over the weights, every realized gate density
satisfies $\alpha_\ell(x)\in[c_0,c_1]$ for all inputs $x$ on the sphere of interest and all
$1\le\ell\le L$.
\end{assumption}

For ReLU with zero bias and a symmetric input law, $\alpha_\ell(x)\to1/2$ and the assumption holds
with any $c_0<1/2<c_1$; for a bias $b_\ell=\beta_\ell\sqrt{n_\ell}\,\mathbf 1$ the density is
$\Phi(-\beta_\ell/\varsigma_\ell)$ with $\varsigma_\ell$ the preactivation scale, again in $(0,1)$; the
assumption fails only if a layer is driven fully on or fully off, which is a degenerate initialization.
All constants below depend on $c_0,c_1$ through $\kappa(c_1)$ and $\kcheck(c_0)$.

\subsection{Counting patterns: the arrangement bound}
The estimates we need must hold not for one input but for all inputs at once, and the gates change
with the input. The key combinatorial fact is that, although there are exponentially many gate
patterns in the abstract, the number \emph{realized} by inputs from a $d$-dimensional domain is only
polynomial in the widths for fixed depth.

\begin{lemma}[Realized pattern count]\label{lem:arrangement}
Fix the weights and biases. As $x$ ranges over $\R^d$, the map $x\mapsto(D_1(x),\dots,D_L(x))$ takes at
most $(C\,\bar n)^{\,d L}$ distinct values, where $\bar n=\max_\ell n_\ell$ and $C$ is absolute. For a
piecewise-linear activation with $r$ slopes the bound is $(C r\bar n)^{dL}$.
\end{lemma}

\begin{proof}
Each neuron's slope changes only across the zero set of its preactivation, which for layer $1$ is a
hyperplane in $\R^d$. The first-layer gates partition $\R^d$ into at most $\binom{\le n_1}{\le d}$-many
$\le(Cn_1)^d$ cells (an arrangement of $n_1$ hyperplanes in $\R^d$ has $O(n_1^d)$ cells; for $r$
slopes each neuron contributes $r-1$ parallel hyperplanes, giving $(Crn_1)^d$). On each first-layer
cell $h_1$ is affine in $x$, so the second-layer preactivations are affine and their zero sets are
again hyperplanes in $x$, refining each cell into $\le(Cn_2)^d$ subcells, and so on. After $L$ layers
the total number of cells, on each of which the whole gate tuple is constant, is at most
$\prod_{\ell=1}^L(Cn_\ell)^d\le(C\bar n)^{dL}$.
\end{proof}

The logarithm of this count is $O(dL\log\bar n)$, which is what we require to be $o(\min_\ell n_\ell)$
for the uniform estimates of the next section; this is the meaning of the \emph{wide regime},
$dL\log\bar n=o(\min_\ell n_\ell)$. It is satisfied whenever the widths exceed the input dimension by a
factor growing faster than $\log$, at fixed depth --- the regime of \cite{bartlettBC21} and the one in
which random networks are provably fragile.
